% !TeX root = ../main.tex

\begin{problem}{Arfken 1-23}
    Problem goes here.
\end{problem}

\begin{solution}
    Write your solution here.

    \section{Equation Labelling and Referencing}
    Equations are automatically numbered. Label equations using \verb|\label| and reference using \verb|\cref|. For example:
    \begin{equation} \label{eq:force}
        F = ma
    \end{equation}
    Based on \cref{eq:force}, force is directly proportional to mass.

    \section{Multiline Equations}
    Create multiline equations using \verb|\begin{align} ... \end{align}|:
    \begin{align}
        y &= mx + b, \\ 
        x &= \frac{y-b}{m}.
    \end{align}
    The equations equations align on \verb|&| line break on \verb|\\|.

    \section{Matrices}
    Create matrices using \verb|\begin{matrix} ... \end{matrix}|. 
    
    (You may also use \texttt{vmatrix, Vmatrix, pmatrix, Pmatrix, instead of matrix}).
    \begin{equation}
    \begin{matrix}
        1 & 2 & 3 \\
        4 & 5 & 6 \\
        7 & 8 & 9
    \end{matrix}
    \to
    \begin{pmatrix}
        \alpha & \beta & \gamma \\
        \theta & \phi & \epsilon \\
        \psi & \omega & \pi
    \end{pmatrix}
    \end{equation}
\end{solution}
